\documentclass[a4paper,11pt,oneside]{sbook}
\linespread{1.1}
%\input{commands}
%% Development and implementation of real space poisson solver and density functional method for investigation finite sized confined systems
%% %% \usepackage{mathrsfs}
%% \usepackage{mathabx}
%\usepackage{graphicx}
\usepackage{pifont}
\usepackage{latexsym,amsfonts,amssymb,mathrsfs,fancyhdr}
\usepackage{amsmath}
\usepackage{makeidx}
%\usepackage{pstricks}
\usepackage {pstcol, pst-grad, epsfig}
\usepackage[dvipdfm,hyperindex,pdfmark,pdftitle={MSC REPORT DRAFT COPY},pdfauthor={SANDIP DE},bookmarks=true,colorlinks=true,nesting]{hyperref}
\usepackage{colortbl}
%\usepackage{eso-pic}	% For background image
\usepackage{color}
\usepackage{eepic}
\usepackage{pifont} 		%This allowlos fancy "Digbats". Refer to:  ~/manuals/tex/symbols-a4.pdf
 \usepackage{yfonts}
\usepackage{textcomp,bbm}
\usepackage{sidecap}
\usepackage{graphicx}
\usepackage{subfig}
%\usepackage{pst-pdf}
%\usepackage[]{hyperref}
%\usepackage{pstricks,pst-3dplot}
\usepackage{listings}
\usepackage{float}
\usepackage{a4wide}
\usepackage{fancyhdr}
\usepackage{amssymb}
\usepackage[newcentury]{quotchap}
\usepackage[Bjornstrup]{fncychap}
%% \usepackage{algorithmic}
%% \pagestyle{fancyplain}
\setlength{\headheight}{12 pt}
\pagestyle{fancy}
\renewcommand{\chaptermark}[1]{\markboth{#1}{}}
\renewcommand{\sectionmark}[1]{\markright{#1}{}}
\fancyhf{}
\fancyhead[LE,RO]{\thepage}
\fancyhead[RE]{\textit{\nouppercase{\leftmark}}}
\fancyhead[LO]{\textit{\nouppercase{\rightmark}}}
\fancypagestyle{plain}{ %
\fancyhf{} % remove everything
\renewcommand{\headrulewidth}{0pt} % remove lines as well
\renewcommand{\footrulewidth}{0pt}}
\newcommand{\HRule}{\rule{\linewidth}{0.8mm}}
\begin{document}
%\newfont {\santitle}    	{pagdo8z at 8 mm} 
\begin{center}
% Upper part of the page
\includegraphics[width=0.3\textwidth]{logo}\\[1cm]
\textsc{\Huge {Advanced Computation Lab \\[0.3cm]
					Department of Physics \\[0.3cm]
					 University of Pune}}\\[2.5cm]
\definecolor{titlecolor}{rgb}{0.0,0.4,0.8}
\definecolor{black}{rgb}{0,0,0}
% Title
\HRule \\[1cm]
\linespread {2}
%\santitle {  Solving Schr\"odinger equation and Poisson's equation for finite sized system and use them to build three dimensional
%  real-space Density Functional Theory code for Confined electron system}\\[0.4cm]
\color{titlecolor}{ \bfseries { \sanlarge {The User Manual \\ [1cm] of The 3d DFT Program }}} \color{black} \\[1cm]
\HRule  \\[1.5cm]
\sanlarge{Sandip De} \\
\today 
\thispagestyle{plain}
%% \end{titlepage}
\end{center}
\chapter*{Understanding The Program}
 This document deals with the technical details of the 3d-dft program. If you have no idea about , what this program does, read my M.Sc project report .
 The whole program can be used using graphical interface or manually executing the program file. 
\section{Structure Of The Program}
 In the 3d-dft directory you will see several files and directories.
 \begin{itemize}
 \item Src : This directory contains the source codes of the programs. Do not change these files unless you know , what you are doing.
 \item bin : The executable file is kept here after compilation of the source code.
 \item input: Here the input configuration files are kept.
 \item info: All output informations while running the program are kept here.
 \item output: The output charge density , input potentials etc can be found here.
 \item charge\_density\_for\_vmd : The charge densities are written in the format required by vmd for visualization.
 \item output\_for\_vmd: wave functions are written in the format required by vmd for visualization.
 \item backup: Here the backup of the outputs are created after predefined number of iterations and the backup files can be used to resume a unfinished job.
 \end{itemize} 
 All other files and directories are required to run the program correctly and are not for your use. 
 \section{Compiling The Program}
  For compiling the program you need to change the makefile in the \emph{src} directory by hand.
  You need to give the exact \emph{MKL Library path}. For example if the makefile contains a line  \newline
  \emph{LFLAGS = -L /opt/intel/mkl/9.1.023/lib/32/ } \newline
  where  \emph{ /opt/intel/mkl/9.1.023/lib/32/} is the MKL library path for my system , you need to change the path by your MKL library path.
  This are very standard practice in program compilations . If you are not familiar with this thing, please contact someone in your lab. You may have to 
  change additional parameters depending upon your system.You need to have fftw3.2 installed and the Library should be in your path.
  do  \newline \emph  {cd src \newline make} \newline to compile the program. After successful compilation the executable file \emph{3d\_dft.x} will be placed in the \emph{bin} directory.
  \subsection{Cleaning after compilation}
  This is not a necessary step.
         You can clean by \emph{make clean}. DON'T DELETE ALL THE *.mod FILES BY HAND. THERE ARE SOME ADDITIONAL mod FILES REQUIRED BY THE COMPILATION.
\section{RUNNING THE PROGRAM}
 If you have successfully compiled the program, you can start using the program by executing the \emph{3d\_dft.x} file. Several scripts are written to make your work easy.
 run the \emph{start.sh} for graphical interface. You will need  
 \begin{itemize}
 \item zenity 	Version 2.24.0
\item  dialog  Version	1.1-20080316    
 \end{itemize} 
 for running the scripts properly. The graphical menu is self explainatory.
 For any help contact 1sandipde@gmail.com. 
\end{document}